\section{Evaluation Protocol}
\label{sec:setup}

\subsection{Cases and Methods}
\label{sec:setup-cases}
The 68 cases combine the two Kusiak--Song wind data sets~\cite{Kusiak2010,Bansal2017} (Table~\ref{S-tab:S-wind}) with circular farms of radii $r=500$, 750 and 1000~m and $N=2,\ldots,10$, $2,\ldots,12$ and $2,\ldots,15$ turbines ($R=38.5$~m, $\ell_{\min}=4D=308$~m, $C_T=0.8$, $K=0.075$).

The \emph{main comparison} includes eight methods ($N_p=30$ where applicable): PSO-VNS (Algorithm~\ref{alg:hybrid}, $\omega=0.5$); PSO; SSA-VNS; SSA; LX-SSA; DE/rand/1/bin~\cite{Storn1997} (scale factor 0.5, crossover rate 0.9); VNS; and MS-SLSQP, a multistart sequential least-squares programming method (SciPy's SLSQP~\cite{Kraft1988}, forward-difference gradients) that minimizes the wake loss under explicit constraints, restarting from the initial, then random layouts (Section~\ref{S-sec:S-protocol}); all others minimize $F_p$~\eqref{eq:penalty} over $[-r,r]^{2N}$. The \emph{component analysis} (Section~\ref{sec:ablation}) adds LX-SSA-VNS and two random-sampling controls (Section~\ref{sec:template}; for RSD-VNS in other studies see Table~\ref{tab:design}): after the common initial population, \emph{RSD-VNS} samples its Phase-1 layouts uniformly in the farm disc and \emph{RS-VNS} in the bounding square (3,015 layouts versus the hybrids' 3,030 Phase-1 evaluations). A split study varies $\omega$ on twelve cases (Section~\ref{sec:split}). Table~\ref{tab:design} lists all studies; the re-evaluations and diagnostics use stored layouts and runs.

\emph{Choice of the pool.} The pool contains what the research questions need: the PSO baseline in two settings, the salp-swarm methods and their hybrids, the local search alone (VNS), a gradient-based method and the random-sampling controls. In the main comparison, DE/rand/1/bin represents the real-coded evolutionary family, whose operators act directly on the continuous coordinates. In the additional experiments, designed after the main comparison, a standard real-coded GA was run on the 68 cases, and on the Horns Rev~1 block and IEA37 Case Study~1 at 6,030 and 30,030 evaluations ($N_p=30$; binary tournament, SBX~\cite{Deb1995} with $\eta_c=15$ and $p_c=0.9$, polynomial mutation~\cite{DebGoyal1996} with $\eta_m=20$ and $p_m=1/n$ for $n=2N$ variables, clipping to the box and elitism, all untuned); it enters the nine-method benchmark pool of Table~\ref{tab:ga-main} and the site pools of Tables~\ref{tab:hr-site}--\ref{tab:iea-means} (details: Section~\ref{S-sec:S-xp-ga}), not the eight-method analysis. The GAs of~\cite{Mosetti1994,Grady2005} optimize the occupation of grid cells, a discrete formulation that cannot be compared directly on continuous coordinates. ACO~\cite{Eroglu2012}, BBO~\cite{Bansal2017} and adaptive GAs were not re-implemented, to bound the number of runs (Table~\ref{tab:design}); their published results use other budgets and constraint handling (Section~\ref{sec:literature}). The pool is thus a controlled set of baselines and controls, not a survey of WFLOP optimizers (Section~\ref{sec:limits}).
% CHECK-FINAL [C61]: RS-VNS and RSD-VNS switch after round(0.5 x 6,030) = 3,015 Phase-1 evaluations; RSD-VNS = experiment rsdisc (68 cases x 30 seeds only; not in the budget, feasible-start, Horns Rev or IEA37 studies); CHECK-DIAG [D17]: RS-VNS Phase-1 replay

\begin{table}[!htbp]
\centering
\caption{Experimental design. Budget: objective evaluations per run, all counted (initial population, finite-difference evaluations of MS-SLSQP); r: random, f: feasibility-preserving initialization. Runs: optimization runs made for that study (``reused'': taken from another study in this table). Ten methods: the eight of the main comparison, LX-SSA-VNS and RS-VNS. Split cases: middle and largest $N$ per farm and data set; six largest: largest $N$ per farm and data set. In total, 52{,}930 optimization runs (16{,}740 in the additional experiments and sensitivity studies).}
\label{tab:design}
\scriptsize
\setlength{\tabcolsep}{3pt}
\begin{tabular}{@{}>{\raggedright\arraybackslash}p{2.0cm}>{\raggedright\arraybackslash}p{1.85cm}>{\raggedright\arraybackslash}p{5.6cm}>{\raggedright\arraybackslash}p{2.15cm}>{\raggedright\arraybackslash}p{1.05cm}>{\raggedright\arraybackslash}p{1.15cm}>{\raggedright\arraybackslash}p{0.75cm}@{}}
\toprule
Study & Cases & Methods & Budget (init.) & Seeds & Runs & Sec. \\
\midrule
Main comparison & 68 benchmark & PSO-VNS, PSO, SSA-VNS, SSA, LX-SSA, DE, VNS, MS-SLSQP & 6,030 (r) & 1--30 & 16{,}320 & \ref{sec:overall} \\
Component analysis & 68 & LX-SSA-VNS, RS-VNS, RSD-VNS (2{,}040 each); six reused & 6,030 (r) & 1--30 & 6{,}120 & \ref{sec:ablation} \\
Old PSO setting & 68 & PSO, $w=0.7$, $c_1=c_2=2$ & 6,030 (r) & 1--30 & 2{,}040 & \ref{sec:psosetting} \\
Budget split & 12 split cases & PSO-VNS, $\omega=0.25$, 0.75, 0.9 (360 each); $\omega=0.5$, 1 reused & 6,030 (r) & 1--30 & 1{,}080 & \ref{sec:split} \\
Coefficient sweep & 12 split cases & PSO, $w=0.7$, $c_1=c_2=c=1.2$, 1.4, 1.6--2.0 (step 0.1); $c=2$ with $v\gets0$ or $\lvert v\rvert\le0.2\,(u-l)$ (360 each) & 6,030 (r) & 1--30 & 3{,}240 & \ref{sec:psosetting} \\
Budget scaling & six largest & ten methods (6,030 reused) & 30,030; 120,030 (r) & 1--30 & 3{,}600 & \ref{sec:budget} \\
Feasible starts & six largest & ten methods except RS-VNS & 6,030 (f) & 1--30 & 1{,}620 & \ref{sec:feasinit} \\
Horns Rev~1 & 16 turbines & ten methods (f: except RS-VNS) & 6,030 (r, f); 30,030, 120,030 (r) & 1--30; 1--10 at 120,030 & 970 & \ref{sec:hornsrev-site} \\
IEA37 CS~1 & 16, 36 turbines & ten methods & 6,030; 30,030 (r) & 1--30 & 1{,}200 & \ref{sec:iea37} \\
\multicolumn{7}{@{}l}{\emph{Additional experiments and sensitivity studies}} \\
Spacing $5D$, $6D$ & 12 split cases & the eight of the main comparison and RSD-VNS ($4D$: stored runs reused) & 6,030 (r) & 1--30 & 6{,}480 & \ref{sec:spacing} \\
Laplace ablation & 12 split cases & LX-SSA without re-evaluation; Laplace candidate replacing the midpoint; $\gamma\sim U(-2,2)$; $\chi=0.25$, 0.5, 2 (360 each); SSA, LX-SSA reused & 6,030 (r) & 1--30 & 2{,}160 & \ref{S-sec:S-xp-lap} \\
Real-coded GA & 68; Horns Rev~1; IEA37 CS~1 & GA (other methods reused) & 6,030 (r); sites also 30,030 (r) & 1--30 & 2{,}220 & \ref{S-sec:S-xp-ga} \\
Exact gradients & IEA37 CS~1, 16 and 36 turbines & MS-SLSQP and PSO-SLSQP with exact gradients & 6,030; 30,030 (r) & 1--30 & 240 & \ref{S-sec:S-xp-grad} \\
Lillgrund & 16 turbines (block) & the eight of the main comparison and RSD-VNS & 6,030; 30,030 (r) & 1--30 & 540 & \ref{sec:lillgrund} \\
Constraint handling & 6 (3 per data set) & PSO-VNS, PSO, GA, SSA-VNS, DE; three constraint handlings (900 each) & 6,030 (r) & 31--60 & 2{,}700 & \ref{sec:boundary} \\
Direct 1$^\circ$ optimization & 8 (4 per data set) & PSO-VNS, PSO, GA, MS-SLSQP, RSD-VNS; 1$^\circ$ objective and seed-paired 15$^\circ$ control (1{,}200 each) & 6,030 (r) & 31--60 & 2{,}400 & \ref{sec:direction} \\
\midrule
Re-evaluation & stored final layouts & cubic power curve, Gaussian wake (main comparison); 5$^\circ$ and 1$^\circ$ bins: 21{,}716 layouts (11 methods), 1{,}030 Horns Rev~1 and 540 Lillgrund runs (also PyWake NOJ) & -- & -- & -- & \ref{sec:robustness} \\
Diagnostics & 6; 12; 7 & instrumented: PSO (both settings); Phase~2 of PSO-VNS, SSA-VNS, RS-VNS; PSO, DE from feasible starts & as stored & 10; 10; 3 & 522 (reruns) & \ref{S-sec:S-diag} \\
\bottomrule
\end{tabular}
\end{table}
% CHECK (manual, 2026-09-29; run counts derived from the design constants and confirmed by the row counts of the run files in analysis/, no macro unless named):
%   main 8 x 68 x 30 = 16,320 = \NRunsMain; component analysis 3 x 68 x 30 = 6,120 new (LXBV in fresh_bgrid, mpce_rsvns, mpce_rsdisc = \NDRsDiscRuns), \NRunsBenchmark = 16,320 + 6,120 = 11 x 2,040;
%   split: mpce_psosplit 720 rows (omega 0.25 / 0.75) + mpce_omega90 360 rows (\NDOmegaNinetyRuns) = 1,080; csweep: mpce_csweep_s0of1.csv 3,240 rows = 9 settings x \NSRuns;
%   budget: mpce_b30k + b30kp 2,100 rows and mpce_b120k + b120kp 1,900 rows, minus the superseded Horns Rev rows (300 and 100; replaced by hrfix) = 1,800 + 1,800 = 3,600 (10 x 6 x 30 x 2);
%   feasible init.: mpce_feas 1,680 + mpce_feasp 210 = 1,890 rows minus 270 superseded Horns Rev rows = 1,620 (9 x 6 x 30);
%   Horns Rev: mpce_hrfix 970 rows = 10 x 30 (6,030) + 10 x 30 (30,030) + 10 x 10 (120,030) + 9 x 30 (feasible) = \NFHRNRuns;
%   IEA37: mpce_iea16/iea36(+p) 1,200 rows = 10 methods x 2 scenarios x 2 budgets x 30 (Table S-tab:iea37-detail lists all ten; the main-text Table tab:iea37 shows eight);
%   total 22,440 + 2,040 + 1,080 + 3,240 + 3,600 + 1,620 + 970 + 1,200 = 36,190;
%   additional experiments (rev2_*_s*of*.csv): 6,480 + 2,160 + 2,220 + 240 + 540 = 11,640 (47,830 with the above; caption total verified 2026-10-04);
%   sensitivity studies (2026-10-04): rev3_constraint_pen/deb/proj.csv 3 x 900 rows = 2,700 (5 methods x 6 cases x 30 seeds per variant); rev3_fine_s0..9of10.csv 1,200 rows
%   + rev3_fine15_s0of1.csv 1,200 (8 cases x 5 methods x 30 seeds per arm; control arm complete, 1,200 rows) = 2,400;
%   grand total 36,190 + 11,640 + 2,700 + 2,400 = 52,930 (16,740 in the additional experiments and sensitivity studies, designed after the main comparison); re-evaluation: rev3_sites.json hr_eval.n_layouts 1,030 (970 + 60 GA), lg_eval.n_layouts 540;
%   diagnostics \NDVerRuns = 2 x \NDDynRuns + 3 x \NDVnsRunsPSOVNS + \NDFsRunsPSO + \NDFsRunsDE (120 + 360 + 42); CHECK-DIAG [D01], [D17]

\subsection{Comparison Design and Statistics}
\label{sec:setup-stats}
\emph{Equal budgets and paired seeds.} Unless stated otherwise, a run uses exactly $B=6{,}030$ objective evaluations. The budget is a model-query budget: it counts objective calls, including the initial population, finite-difference calls and repeated evaluations of a layout, and charges an exact gradient (IEA37 only) as $c_g=3$ objective equivalents (Section~\ref{sec:iea37}). Optimizer overhead, constraint evaluations, SLSQP subproblems and the geometric packing of the feasibility-preserving initialization are not charged, so equal budgets do not imply equal elapsed times (Section~\ref{sec:runtime}). Each method uses seeds 1--30 and draws its initial population first, so runs with equal seeds share it and are paired.
% verified 2026-09-28: first Curve value identical across fresh_grid/vgrid/bgrid and mpce_psoc/psobv/rsvns/slsqp (13 files) for all 792 case-seed pairs with a feasible initial best (else NaN)

\emph{Feasibility-aware ranking.} Let $f_{a,c}$ be the number of feasible runs of method $a$ in case $c$ (tolerance $10^{-6}$~m; Section~\ref{S-sec:S-tol}) and $\bar L_{a,c}$ the mean wake loss of these runs. Method $a$ is \emph{qualified} in case $c$ if at least half of its runs are feasible (15 of 30, 5 of 10 for 10 seeds). Within a case, the qualified methods are ranked by the mean benchmark objective of their feasible runs (rounded to $10^{-6}$, so that exact ties, e.g., at the wake-free optimum for $N=2$, share their average rank), and the others below them, by their number of feasible runs. This qualification-then-conditional-quality rule is not a feasibility-first ordering of all runs: once two methods qualify, 15/30 feasible runs can outrank 30/30 if their conditional mean is better, so it limits, but does not remove, survivorship bias. The ranks therefore summarize conditional solution quality subject to a minimum success threshold, not reliability; we report success counts separately and add an all-run paired outcome for the 68 cases, the Lillgrund block (Section~\ref{sec:archive-reliability}) and the spacing study (Section~\ref{sec:spacing}).

\emph{Inference hierarchy.} Claims that two methods differ rest on the case-mean Wilcoxon signed-rank tests (item~2 below), the primary test of the main comparison, Holm-corrected within each table; as the 68 cases are designed, not sampled, their $p$ values describe this benchmark, not a population of farms. Claims of practical equivalence or non-superiority, among them the two primary questions below, are decided at the seed level, for the fixed benchmark (Definition~\ref{def:S-equiv}(a)); the case and cluster levels, joint resampling of the synchronized seeds and equal-cluster weighting are sensitivity analyses. Average ranks are descriptive summaries that depend on the method pool; their Friedman and post hoc tests are reported for comparison with common practice, and the run-level W/T/L counts are tallies. The all-run paired outcome is a separate, secondary endpoint: it counts wins and does not measure the size of a difference. The additional experiments and sensitivity studies (Table~\ref{tab:design}) were designed after the main comparison and are secondary to it.

\emph{Tests.} All tests are two-sided ($\alpha=0.05$), with Holm's correction~\cite{Holm1979} within each family of comparisons.
1)~\emph{Case ranks}: the Friedman test on the case-by-method rank matrix with the Iman--Davenport $F$ statistic, and post hoc $z$ tests of the average rank $\bar R_a$ of every method against that of PSO-VNS, $z=(\bar R_a-\bar R_{\text{PSO-VNS}})/\sqrt{k(k+1)/(6n)}$ for $k$ methods and $n$ cases~\cite{Demsar2006,Derrac2011,LaTorre2021}.
2)~\emph{Case means} (the primary test for differences, as mean-rank post hoc tests depend on the pool~\cite{Benavoli2016}): for methods $A$ and $B$, the case-mean difference is $d_c=\bar L_{A,c}-\bar L_{B,c}$ (percentage points, pp; negative: $A$ better), and the Wilcoxon signed-rank test is applied to the $d_c$ of each compared pair (in the main comparison, PSO-VNS against each method); a case in which only one of the two methods is qualified counts as a maximal difference in its favor (larger than every observed difference; cases: 2 for SSA, 2 for LX-SSA and 18 for DE), and one in which neither is qualified is dropped (0 cases). With such cases, the signed-rank test concerns this constructed combined score, which mixes wake-loss differences with qualification outcomes, rather than wake loss alone; $\overline{\Delta L}$ excludes them. The mean $\overline{\Delta L}$ of the $d_c$ over the cases in which both are qualified carries a 95\% percentile bootstrap CI (10{,}000 resamples of the cases, fixed seed). For key pairs we also give the Hodges--Lehmann estimate that belongs to the Wilcoxon test (median of the Walsh averages $(d_c+d_{c'})/2$, $c\le c'$; 95\% CI by exact test inversion; Table~\ref{S-tab:X-hl}), which can differ from $\overline{\Delta L}$ when a few large differences dominate the mean.
3)~\emph{Runs}: per case, Wilcoxon signed-rank tests on the 30 seed-paired runs, scored feasibility first (infeasible runs below every feasible run and ordered by minimum spacing; differences $\lvert d\rvert\le10^{-9}$ dropped), with the rank-biserial correlation $r_{\rm rb}$ and the counts of cases with a significant win, no significant difference and a significant loss (W/T/L, descriptive tallies).
\emph{Holm families.} A family is the set of comparisons of one table: for the rank post hoc tests, the $k-1$ comparisons with PSO-VNS of one Friedman analysis; for the case means, the comparisons of PSO-VNS with the other seven methods of the main comparison, or the fifteen contrasts of the component analysis (Section~\ref{sec:ablation}); for the run-level tests, the comparisons of one table within one case. The main-comparison and the component-analysis tables are different families, so a pair can have slightly different W/T/L tallies in the two; one family over all 28 case-mean tests is examined below.
% CHECK (manual, 2026-09-29): constants in analysis/mpce_results.py: BOOT_N = 10000 (fixed seed), wil() drops |d| <= 1e-9, rank_rule rounds to 1e-6 and ranks non-qualified methods by nfeas (average ranks for ties); Holm families: friedman_block (vs focus), case_mean_wilcoxon (per table), paired_vs / ablation contrasts (per case and table)

\emph{Practical equivalence.} A nonsignificant test does not show that two methods are equally good, and a significant difference can be negligible; we therefore assess equivalence by interval inclusion, motivated by two one-sided tests (TOST)~\cite{Lakens2017}. \emph{Primary questions.} Two equivalence questions are primary: whether the VNS phase adds a practically relevant gain to PSO (PSO-VNS vs.\ PSO, two-sided equivalence) and whether an SSA first phase gives a practically relevant gain over disc sampling (SSA-VNS vs.\ RSD-VNS, non-superiority). Each is assessed on its own at $\alpha=0.05$; all other equivalence results (further pairs, subgroups, spacings, direction resolutions and the GA comparisons) are exploratory and are not combined into a family-wide equivalence claim.

\begin{definition}[Practical equivalence at three levels of inference]
\label{def:S-equiv}
Let $A$ and $B$ be two methods, $d_c$ their case-mean difference over the $n$ cases in which both are qualified, $\overline{\Delta L}=\frac1n\sum_c d_c$, and $m>0$ a margin. $A$ and $B$ are \emph{practically equivalent at margin $m$} at a given level if both one-sided hypotheses $H_0^-\colon\Delta\le-m$ and $H_0^+\colon\Delta\ge m$ about the estimand $\Delta$ of that level are rejected at $\alpha=0.05$, i.e., if the 90\% interval for $\Delta$ lies inside $(-m,m)$. The three levels are:
\begin{enumerate}
\item[(a)] \emph{Seed level} (the fixed benchmark). Estimand: the benchmark average of the expected case-mean differences; only the run-to-run variability is random. Interval: percentile bootstrap that resamples the 30 seed-paired runs within every case (the same seeds for both methods; case means over the feasible resampled runs).
\item[(b)] \emph{Case level} (exchangeability sensitivity). Estimand: the mean case difference under a hypothetical exchangeable distribution represented by the 68 designed cases. Interval: percentile bootstrap over cases, treating them as independent; a sensitivity analysis, as the cases are a designed benchmark, not a random sample of farms.
\item[(c)] \emph{Cluster level} (benchmark-group sensitivity). Estimand: the case-weighted mean difference allowing dependence within the six designed (data set, radius) groups with nested $N$. Interval: cluster-robust $t$ interval with the bias-reduced CR2 variance and a $t$ distribution with $G-1=5$ degrees of freedom ($G=6$ clusters), and a restricted wild-cluster bootstrap-$t$ TOST (CR2-studentized, all $6^6$ draws of Webb's six-point weights enumerated); equivalence at this level requires both.
\end{enumerate}
The \emph{minimal margin} of a level, $m_{\min}=\max(\lvert\ell\rvert,\lvert u\rvert)$ for the 90\% interval $[\ell,u]$, is the smallest margin at which equivalence holds: it holds for every $m>m_{\min}$ and for none smaller.
\end{definition}
For one pair, the intersection--union principle controls TOST at level $\alpha$ when its constituent tests are calibrated; it does not itself calibrate percentile-bootstrap tail shares or adjust comparisons of several pairs (Section~\ref{S-sec:S-th-equiv}). Every equivalence and non-superiority verdict is decided by interval inclusion (seed and case level: 90\% percentile bootstrap intervals; cluster level: the CR2 interval together with the wild-cluster bootstrap TOST); the bootstrap $p_{\rm TOST}$ agrees with this rule for all 18 pairs at every level (Table~\ref{S-tab:S-sa-tostaudit}), which is algebraic consistency, not a check of the error rate. That check is an empirical shifted-null simulation on this dataset: with the differences shifted to a true mean at the margin, the equivalence decision has an empirical size of 4.4--7.4\% for PSO-VNS vs.\ PSO (case and seed level) and 4.7--6.0\% for SSA-VNS vs.\ RSD-VNS at the seed level, but 15.0\% for the latter's case-level non-superiority test (one outlying case, $d_c=-1.33$~pp), so the nominal 5\% is not attained exactly. Against this shifted null, the case-level sensitivity results for the two primary questions have calibrated $p$ values of 0.015 and 0.007; these are local diagnostics, not a general guarantee of level-0.05 control, and enter no verdict. The seed-level simulation resamples the seeds independently within each case, so it does not calibrate the joint resampling of the synchronized seed vectors (below).
A separate local check draws the seed vectors jointly (5,000 shifted-null replications; Table~\ref{S-tab:S-sa-dep-bench}): the seed-level equivalence decision then has an empirical size of 6.94\% at $-m$ and 4.30\% at $+m$ for PSO-VNS vs.\ PSO (above the nominal 5\% at $-m$) and 5.42\% and 5.68\% for SSA-VNS vs.\ RSD-VNS; the calibrated $p$ values support the two primary verdicts (PSO-VNS vs.\ PSO equivalent, $p<0.001$ on both sides; no gain of SSA-VNS over RSD-VNS larger than $m$, $p_-<0.001$), whereas the equivalence of SSA-VNS and RSD-VNS is still not established ($p_+=0.058$).
$\overline{\Delta L}$ is the same at all levels; its intervals widen here (not a general nesting guarantee), so an equivalence statement must name its level. Seed $k$ seeds the same random stream in every case (the same initial population, scaled by $r$, for equal $N$); resampling one seed vector jointly for all cases widens the seed-level intervals by up to 1.9$\times$, leaves PSO-VNS vs.\ PSO unchanged and moves SSA-VNS vs.\ RSD-VNS to [$-$0.002, 0.050]: its seed-level equivalence is not robust, its non-superiority is (Table~\ref{S-tab:S-sa-syncseed}). At every level, the set of cases in which both methods qualify is fixed from the original data; at the seed level, case means are taken over the feasible resampled runs of each method, so the two methods can be averaged over different seed subsets (for the two primary pairs, at least 99.6\% of the runs are feasible). Equal case weighting gives the clusters with more cases (9, 11 and 14 per radius) more influence and targets the case average of this benchmark, not a typical new scenario group; weighting the six clusters equally leaves every verdict at $\pm0.05$~pp unchanged (PSO-VNS minus PSO: $-0.014$~pp, 90\% CI [$-$0.063, 0.035], also with wild-cluster bootstrap intervals; Table~\ref{S-tab:S-sa-eqclus}). ``$A$ gives no gain larger than $m$ over $B$'' (non-superiority) needs only one of the two tests: the lower end of the 90\% interval must lie above $-m$. A result that is neither equivalent nor significantly different is inconclusive and supports no statement of the kind ``as good as''. The margin $m=0.05$~pp (4.1 and 2.1~MWh/yr per turbine in Data Sets I and II) is a post hoc convention; the minimal margins $m_{\min}$ let readers apply their own, and a confirmatory practical-equivalence claim would need a margin fixed before a new study.

\emph{Bayesian signed-rank test.} The Bayesian signed-rank test~\cite{Benavoli2017} places a Dirichlet-process prior (strength 0.5, pseudo-observation at 0) on the distribution of the $d_c$ and, for every posterior sample, computes the shares $\theta_A$, $\theta_{\rm rope}$ and $\theta_B$ of the Walsh averages $(d_c+d_{c'})/2$ below $-m$, inside the region of practical equivalence (ROPE) $[-m,m]$ and above $m$. We report the posterior probability that equivalence is the most probable of the three outcomes, the share of the 50{,}000 posterior samples in which $\theta_{\rm rope}$ is the largest (Section~\ref{sec:abl-swarm} and Table~\ref{S-tab:X-bayes}). It is not the probability that $\Delta$ lies within $\pm m$; it concerns the typical pairwise-averaged case difference, so it can disagree with TOST, and it treats the cases as exchangeable.
% CHECK-FINAL [C49]-[C52]: margin_pp == EQ_MARGIN; TOST: 90 % bootstrap CI inside (-m, m); Bayes signed-rank, rope (-m, m); min_margin_pp in tab:equivalence
% CHECK-EXTRA [X38]-[X42]: seed / case / cluster (CR2 t(5), wild-cluster Webb) levels and minimal margins (tab:X-equiv-levels); [X56]: margin = \NXMarginMWhTurbDsI / \NXMarginMWhTurbDsII MWh/yr per turbine; [X48]: Bayes = P(rope most probable); [X49]-[X51]: Hodges-Lehmann (tab:X-hl). Model-shift margin justification (X29/X30) deliberately dropped (D13).
% CHECK (manual, 2026-09-29): analysis/mpce_inference_extra.py: seed level resamples the 30 seed indices within every case (paired), cr2() uses t with G - 1 = 5 d.f., webb_all() enumerates 6^6 = 46,656 Webb draws, cluster-level equivalence = CR2 CI inside +-m and wild p_TOST <= 0.05; mpce_results.py: tost() add-one corrected tail shares, BAYES_S = 0.5, BAYES_Z0 = 0, BAYES_N = 50000. Definition adapted from optA/supp_theory.tex (def:S-equiv moved to the main text, see optA/sw/moved.md)

\emph{Robustness of inference.} Six qualification thresholds (1, 10, 15, 20, 25 and 30 feasible runs), dropping the cases with an imputed maximal difference and keeping zero differences (Pratt; Table~\ref{S-tab:S-sa-imputation}) change no verdict; cluster analyses (sign per cluster, cluster-robust $t$ with CR1 variance and 5 degrees of freedom, cluster bootstrap, leave-one-out) qualify only the significant contrasts of the salp-swarm hybrids with the controls; one Holm family over all 28 case-mean tests makes only PSO-VNS vs.\ PSO on the cases with $N\ge 10$ (both data sets) and PSO-VNS ($\omega =0.5$) vs.\ PSO-VNS ($\omega =0.9$) nonsignificant (Tables~\ref{S-tab:X-threshold}--\ref{S-tab:X-multiplicity}).
% CHECK-EXTRA [X02]-[X09]/[X31]: six thresholds (NXThrList) change no verdict; [X10]: six clusters; [X11]: significant pairs unanimous in 6/6 clusters except \NXClCaseSigNotUnanimous (the significant contrasts of the salp-swarm hybrids with the controls are exactly these two plus SSA-VNS vs RS-VNS; LX-SSA-VNS vs RS-VNS is not significant, [X13]); [X16]/[X17]: cluster-robust t not significant for the two RSD-VNS contrasts and SSA-VNS vs RS-VNS (p \NXClSSAVNSvsRSVNSCRP); [X14]: LOCO changes only \NXLocoChanged; [X18]: one Holm family over the \NXMultN case-mean Wilcoxon tests loses \NXMultLostHolm. CHECK-DIR [F01]: 1-deg sub-bin re-evaluation (one sub-bin reproduces the benchmark objective)
% CHECK (manual, 2026-09-29): tab:X-cluster caption (mpce_supp_inference.tex): CR = cluster-robust t of the case-weighted mean, CR1 variance, 5 d.f.

\subsection{Tests Beyond the Benchmark}
\label{sec:setup-beyond}
\emph{Initialization and budget.} The feasible initialization draws seeded layouts inside the site and separates the turbines with L-BFGS-B on a packing penalty, without objective evaluations (Section~\ref{S-sec:S-protocol}); all methods but RS-VNS use it at 6,030 evaluations. With random starts, all ten run at 6,030, 30,030 and 120,030 evaluations. Both studies use the six largest cases and the Horns Rev~1 block, with 30 seeds (10 for Horns Rev~1 at 120,030).

\emph{Other sites and models.} The Horns Rev~1 block is PyWake's 16-turbine example~\cite{PyWake}: V80 turbines ($D=80$~m), the measured 12-sector wind climate, the parallelogram through the corner turbines as boundary (Section~\ref{sec:hornsrev-site}), $\ell_{\min}=4D$ and the Jensen wake with $K=0.04$ and $C_T$ at the free-stream speed. A block of the Lillgrund offshore wind farm is a second site with a measured wind climate (Section~\ref{sec:lillgrund}).
% CHECK (manual, resolved 2026-09-29): description matches analysis/hornsrev_model.py (V80, 12 sectors, K = 0.04, hub-centre top-hat, C_T at free-stream speed, 5-deg bins after the hrfix) and supplement sec:S-hrmodel; all HR runs are the hrfix reruns
The 16- and 36-turbine scenarios of IEA37 Case Study~1~\cite{Baker2019,IEA37repo} (circular sites, radii 1300 and 2000~m, $\ell_{\min}=2D$, 16 directions at one speed, a simplified Bastankhah Gaussian wake) test optimizers rather than represent a site; the ten methods of the budget study, GA and two exact-gradient SLSQP variants (13 methods) run on them at 6,030 and 30,030 evaluations (30 seeds; Section~\ref{sec:iea37}).
% CHECK (manual, 2026-09-29): mpce_iea16/iea36(+p) contain PSOBV, PSOC, SSABV, SSA, LXSSA, DE, BVNS, SLSQP, LXBV, RSVNS (120 runs each); formerly "the eight methods of the main comparison run on them"; 2026-10-04: + GA (rev2_gaiea) and MS-SLSQP / PSO-SLSQP with exact gradients (rev2_grad) = 13-method pool of analysis/rev3_sites.json iea_pool

\subsection{Verification and Reproducibility}
\label{sec:validation}
\emph{Generated results and automated checks.} Every table, figure and number of the article and the supplement is generated from the stored per-run records by the analysis scripts, and every outcome-dependent statement is linked to an automated condition. For validation, the complete analysis pipeline was re-executed from the stored records in the reproducibility package (main study, additional experiments and sensitivity studies): all 208 regenerated or compared tables, number macros, summaries and figures, 117 of them from the sensitivity analyses, are identical to the stored pipeline outputs apart from time stamps, recorded run times and software versions; all 62 main-result, 57 statistical, 22 diagnostic, 41 direction-resolution, 15 coefficient-sweep checks and the 57 conditions of the theory checks T01--T13 pass, as does the audit of the records of the additional experiments; and their analysis reproduces its summary (4{,}670 values) exactly. The analyses of the sensitivity studies and of the precision reruns were regenerated in the same way (Section~\ref{S-sec:S-archive}). Of the stored optimization runs, the validation reran a determinism sample (four runs of PSO-VNS and SSA-VNS), and the precision audit reran the 5{,}714 records whose feasibility labels the rounded coordinates cannot decide, both with the archived drivers: gradient-free runs reproduce their records bit for bit or to floating-point noise, except for the elapsed time, whereas SLSQP-based runs (MS-SLSQP, exact-gradient SLSQP) are not bit-reproducible across CPU and BLAS builds (Section~\ref{sec:archive-reliability}). Not re-executed were the optimization runs themselves, the instrumented diagnostic reruns (below; their stored outputs pass the diagnostic checks) and the calibration runs of the original code (recomputed from their stored outputs).
% CHECK (manual, 2026-10-04; typed counts, no macro): re-execution from the package, repro_package_build/verification/verification_summary.txt: 91 outputs (58 byte-identical, 33 identical apart from time stamps), check_final 62 PASS / 0 FAIL, check_extra 57/57, check_diag 22/22 (stored diagnostics summary), check_dir 41/41, check_csweep 15/15, theory 57 PASS, rev2_summary.json 4,670 values, audit_archive.py PASS, verify_determinism.py 4/4 bit-identical except Seconds (PSO-VNS, SSA-VNS); C4 reruns (analysis/rev3_precision_rerun_summary.json, 2026-10-04 22:11): all diverged reruns are SLSQP-based (MS-SLSQP of fresh_hr16, mpce_iea, rev2_grad). Earlier (2026-09-29): analysis/check_final.log "62 PASS, 0 FAIL, 0 PENDING, 0 TAG ERRORS"; check_extra.log "57 PASS / 0 FAIL"; check_diag.log "22 PASS / 0 FAIL / 0 PENDING"; check_dir.log "41/41 checks passed"; check_csweep.log "15 PASS / 0 FAIL / 0 PENDING"; make_theory_figures.py --check "57 PASS, 0 FAIL" (optA/reviews2/R5_claims_audit.md; re-run on a scratch copy 2026-09-29).

\emph{Instrumented diagnostics.} The diagnostics (Table~\ref{tab:design}) rerun stored runs with instrumented copies of the optimizers, which differ from the code of the paper only by logging lines that read the state and draw no random numbers. All 522 instrumented runs reproduce the stored runs (wake loss, coordinates, convergence curve, evaluations), and seed~1 of every case and method, rerun with the original, uninstrumented code in the same process, is bit-identical (62 of 62 runs). A replay of the random stream of Phase~1 of RS-VNS agrees with the stored runs (120 of 120).
% CHECK-DIAG [D01]: instrumented runs reproduce the stored runs (\NDVerStoredMatch/\NDVerRuns); [D02]: seed-1 reruns with the original classes bit-identical (\NDVerBitMatch/\NDVerBitRuns); [D17]: RS-VNS Phase-1 replay agrees; CHECK-CSWEEP [S01]: c = 2.0 reproduces the 360 stored old-setting runs bit for bit

\emph{Evaluator and code validation.} Re-executed from the package: our objective matches the original implementation (largest difference $8.7\times10^{-11}$ over 720 archived final objectives) and the official IEA37 calculator (relative difference ${<}10^{-11}$), and the Horns Rev~1 and Lillgrund models were compared with PyWake's NOJ model at identical bins (Section~\ref{sec:hornsrev-site}). Recomputed from the stored calibration runs (not re-executed): our reused SSA, LX-SSA, PSO and DE code reproduces the recorded result distributions of the original code (2 of 24 Mann--Whitney tests reject at $\alpha=0.05$; 1.2 expected).
% CHECK (manual, re-run 2026-09-29; no macro): analysis/validate_evaluator.py -> 720 records, max|dObj| 8.73e-11; python3 analysis/iea37_model.py -> rel.err 8.8e-12 / 2.8e-12 / 2.9e-12 (16/36/64); analysis/calibration_vs_recorded.csv (calibrate_authors_code.py; regenerated from the stored calibration outputs by analyze_authors_runs.py in the package verification, the calibration runs themselves not re-executed, verification_summary.txt items 2 and 4; Lillgrund vs PyWake: rev2_site_model.py --pywake) -> 2 of 24 P < 0.05, 0.05 x 24 = 1.2. Lead: move into generated macros + check (R5-U3) if time permits
% CHECK-DIR [F31]: PyWake NOJ comparison at identical bins

\emph{Code, data and environment.} Runs: Linux containers (Intel Xeon; Python~3.11.15, NumPy~2.4.6, SciPy~1.17.1). Each run record gives the method, case, seed, budget, initialization, evaluations, final objective, feasibility, minimum spacing, elapsed seconds, coordinates and convergence curve. The reproducibility package (Data availability; Section~\ref{S-sec:S-archive}) contains all records, models, optimizers and analysis scripts, a build script, a README with the command behind every run file and every table, figure and number, checksums and a Dockerfile; it is provided to the editor and reviewers with this submission and will be deposited publicly with a DOI upon acceptance.
% CHECK (manual, resolved 2026-09-29): versions pinned in analysis/requirements.txt ("from the container that produced the results"); clock speed and core count dropped because the analysis container (Xeon 2.1 GHz) and earlier run containers (2.8 GHz, final_prose.py) differ
% CHECK (manual, 2026-09-29): run-file columns Algorithm, Dataset, Radius, Turbines, Seed, Budget, Init, Calls, Objective, Ideal, WakeLoss, Feasible, MinSpacing, Seconds, Coordinates, Curve (mpce_*.csv); analysis/README_reproduce.md (command per run file); analysis/build_swevo.sh / build_mpce_paper.sh (mpce_results.py + check scripts)
